\section{Addressing the Structural-Avoidance Incentive}
\label{sec:avoidance}

The companion paper \citep{lasser2026poset} identifies a perverse incentive
under the myopic objective: discovering subsumption structure can decrease
$\volL$ (Proposition~5(c)), so the actor is rewarded for breadth over depth.
This section shows that the discounted objective addresses this incentive
under a criterion that holds when the planning horizon is long relative to
the frequency of affected decisions.

\subsection{The Incentive Under the Myopic Objective}

Under the myopic objective $\pi^* = \arg\max \mathbb{E}[\Delta \volL(G) \mid
\pi]$, structural discovery has expected value:
\begin{equation}
\label{eq:myopic_structure}
\mathbb{E}[\Delta \volL \mid \text{discover edge } d_1 \leq d_2]
= \volL(G_{\text{after}}) - \volL(G_{\text{before}})
\leq 0
\end{equation}
by Proposition~5(c) of \citet{lasser2026poset}. The inequality can be strict
when $d_1$ is a shared prerequisite. The myopic actor sees only the immediate
contraction and avoids structural discovery.

\subsection{The Incentive Under the Discounted Objective}

Under $\Vdisc$, the actor evaluates structural discovery by its effect on the
\emph{entire future trajectory}, not just the current step. Structural
knowledge improves the quality of future $\volL$ estimates and planning
decisions. Define the \emph{information value of structure}:

\begin{definition}[Information Value of Structure]
\label{def:info_value}
The \emph{information value} of discovering subsumption edge $d_1 \leq d_2$
    is the difference in expected discounted return between planning with the
    correct poset structure and planning without it:
\begin{equation}
\label{eq:info_value}
I(d_1 \leq d_2) = \mathbb{E}_{\tau \sim \pi^*_{\text{true}}}\!\left[
    \sum_{t=1}^{H} \gamma^t \Delta \volL(G_t)
\right]
- \mathbb{E}_{\tau \sim \pi^*_{\text{false}}}\!\left[
    \sum_{t=1}^{H} \gamma^t \Delta \volL(G_t)
\right]
\end{equation}
where $\pi^*_{\text{true}}$ is the optimal policy under the correct world
    model (with the edge), $\pi^*_{\text{false}}$ is the optimal policy under
    the incorrect model (without it), and each expectation is taken over the
    state trajectory $\tau = (G_1, G_2, \ldots)$ induced by its respective
    policy.
\end{definition}

The information value is the discounted sum of improvements in action quality
that come from having accurate structure. An actor with the wrong model
makes decisions that are systematically suboptimal: it overestimates the
volume of capability sets that share prerequisites (because it doesn't know
about the subsumption) and therefore overinvests in redundant capabilities.

\begin{proposition}[Structural Discovery is Value-Positive Under Discounting]
\label{prop:structure_value}
Under the discounted objective $\Vdisc$, structural discovery is
    value-positive when the information value of the discovered edge exceeds
    the immediate $\volL$-contraction:
\begin{equation}
\label{eq:structure_criterion}
I(d_1 \leq d_2) > |\Delta \volL(G \mid \text{discover } d_1 \leq d_2)|
\end{equation}
\end{proposition}

\begin{proof}
The expected discounted return of each strategy, starting from state $G_0$:
\begin{align*}
\Vdisc(\text{discover}) &= \Delta \volL(G_0 \mid \text{edge}) +
    \mathbb{E}_{\tau \sim \pi^*_{\text{true}}}\!\left[
    \sum_{t=1}^{H} \gamma^t \Delta \volL(G_t)\right] \\
\Vdisc(\text{avoid}) &= 0 +
    \mathbb{E}_{\tau \sim \pi^*_{\text{false}}}\!\left[
    \sum_{t=1}^{H} \gamma^t \Delta \volL(G_t)\right]
\end{align*}
Each expectation is taken over the trajectory distribution induced by its
respective policy, so the two $G_t$ sequences are generally different. The
difference is $\Vdisc(\text{discover}) - \Vdisc(\text{avoid}) =
\Delta \volL(G_0 \mid \text{edge}) + I(d_1 \leq d_2)$. Since
$\Delta \volL(G_0 \mid \text{edge}) \leq 0$, this is positive when
$I(d_1 \leq d_2) > |\Delta \volL(G_0 \mid \text{edge})|$.
\end{proof}

\paragraph{Scope: existence result, not decision procedure.}
Proposition~\ref{prop:structure_value} is an existence result: it identifies
the condition under which structural discovery is value-positive, showing
that the structural-avoidance incentive \emph{can} be overcome. It is not a
decision procedure, because evaluating $I$ requires comparing optimal
policies under two world models, which is itself a planning problem at least
as hard as the one the actor is trying to solve. A computable proxy for $I$
(e.g., an information-gain heuristic or a rollout-based estimate) is needed
for practical use; this is listed as an open problem in
Section~\ref{sec:discussion}.

\paragraph{When the condition holds.} The information value $I$ grows with
the planning horizon $H$ (more future decisions benefit from accurate
structure), the frequency of decisions involving the capabilities $d_1$ and
$d_2$ (more relevant decisions means more improvement), and the degree of
suboptimality of the incorrect model (larger model errors produce larger
action-quality gaps). The immediate contraction $|\Delta \volL|$ is bounded by
the weight of the shared prerequisite. For edges between frequently-used
capabilities with substantial subsumption, the information value
dominates after a modest number of planning steps.

\paragraph{Connection to the compound avoidance loop.} The companion paper
identifies a compound loop: structural avoidance prevents leverage estimation,
which prevents the self-correcting B-to-C dynamic from activating. The
discounted objective breaks the first link in this chain: the actor accepts
structural contraction because the information value is positive over the
planning horizon. Once the structural-avoidance incentive is counteracted,
leverage estimation can proceed (the actor interacts with capabilities it
previously avoided), and the self-correction activates.
